\begin{proof}
   Assume there exists a least $k \in \NN^+$, then it is clear that $\langle g \rangle := \{\dots g^{-2}, g^{-1}, e_G, g, g^2, \dots\} = \{e_G, g, g^2, \dots, g^{k-1}\} $, since $g^k = e_G$ and similarly, $g^{-k} = (g^k)^{-1} = e^{-1}_G = e_G$. Thus, for any $n \in \ZZ$, there exists $q \in \ZZ$ such that
   \[ g^n = g^{qk + \rem(n,k)} = (g^k)^q \cdot g^{\rem(n,k)} = g^{\rem(n,k)}. \] 
   It follows that $|g| = |\langle g \rangle| \leq k$. To see that $|g| = k$, suppose by contradiction that $|g| < k$, then by the pigeonhole principle, there are two elements $0 \leq i < j < k - 1$ such that $g^i = g^j$, hence $g^{j-i} = e_G$ and $j - i \geq 1$ and $j - i < k$, contradicting, $k$ being the least such power. 

   It remains to show that no such $k$ exists if $|g| = \infty$. Indeed, if $|g| < \infty$, then $g^m = e_G$ for some $m \in [1, |g|]$, so if no such $m$ exists, $|g| = \infty$. 
\end{proof}

\begin{defbox}
   A group $G$ is called \textbf{cyclic} if $G = \langle g \rangle$ for some $g \in G$. 
\end{defbox}

\begin{example}
   $(\ZZ, +)$ is a cyclic group with $1$ being a generator, as well as $-1$, since every integer is a multiple of $1$ or of $-1$. 
\end{example}
\begin{example}
   $(\modulo, +)$ is a  cyclic group of order $n$, where $\ol{1}$ is a generator. In fact, if $\gcd(k,n) = 1$, then $\ol{k}$ is also a generator. 
\end{example}

\begin{remark}
   We will show in the future that $\ZZ$, and $\modulo$ are the only cyclic groups up to an ``isomorphism".
\end{remark}

\subsection{Subgroups Generated by a Subset}

It is a fact (proven in the homework) that an intersection of subgroups is a subgroup itself. The proof also works for any arbitrary intersection of subgroups. The following remark follows from this fact:

\begin{remark}
   Let $G$ be a group and $S \subseteq G$, ($S$ can be empty). There exists a smallest (with respect to inclusion) subgroup $H \leq G$ such that $H \supseteq S$ i.e. $H \supseteq S$ and for every $K \leq G$ and $K \supseteq S$, $H \subseteq K$. Namely,
   \[ H := \bigcap_{S \subseteq K \leq G} K. \]
\end{remark}

\begin{defbox}
   Let $G$ be a group and let $S \subseteq G$ ($S$ can be empty). The \textbf{subgroup generated by $S$}, denoted $\langle S \rangle$ or $\langle S \rangle_G$, is the least subgroup containing $S$, i.e. the intersection of all subgroups containing $S$. 
\end{defbox}

\begin{example}
   For a group $G$, $g \in G$, the subgroup generated by $g$ is the same as the subgroup generated by $\{g\}$. 
\end{example}
\begin{example}
   In $(\ZZ, +)$, we have that $\langle \{a,b\} \rangle = \gcd(a,b)\ZZ$ because by Theorem 6.1.2, $\gcd(a,b)$ can be written as a linear combination of $a$ and $b$. 
\end{example}

\subsection{Permutation Groups}

Recall that for a set $X$, we defined the group $S_X$ to be all bijections from $X$ to $X$. These bijections are called \textit{permutations} of $X$ and $S_X$ is called the \textit{group of permutations} or \textit{permutation group} on $X$. In this case, we will mainly focus on permutation groups where $X$ is finite. For these finite $X$, we define $S_n := S_X$ and call $S_n$ the \textit{permutation group of $n$ elements}. We denote the permutations of $S_n$ as $\sigma(1), \sigma(2), \dots$. If $\sigma(i) = i$, we may omit it from the notation. A permutation of $n$ elements can be written like so:

\[ \begin{pmatrix}
1 & 2 & 3 & \dots & n \\
\sigma(1) & \sigma(2) & \sigma(3) & \dots & \sigma(n)
\end{pmatrix} 
\]
where $1 \mapsto \sigma(1)$, $2 \mapsto \sigma(2)$ and so on. Note that since $\sigma$ is a bijection, every element that appears in the top row will appear in the bottom row.

\begin{example}
   Consider this composition in $S_4$:
   \[ \begin{pmatrix}
      1 & 2 & 3 & 4 \\
      2 & 4 & 3 & 1
      \end{pmatrix}
      \circ
      \begin{pmatrix}
      1 & 2 & 3 & 4 \\
      2 & 1 & 3 & 4
      \end{pmatrix}
   \]
   To compute this permutation, we apply the innermost permutation, followed by the outermost one. Let $\sigma \in S_4$ be the outermost permutation and $\tau \in S_4$ be the innermost permutation. We see $\tau(1) = 2$, so $(\sigma \circ \tau)(1) = \sigma(\tau(1)) = \sigma(2) = 4$, so $1 \mapsto 4$. Similarly, $(\sigma \circ \tau)(2) = \sigma(\tau(2)) = \sigma(1) = 2$, so $2 \mapsto 2$, and so on. The composition $\sigma \circ \tau$ is thus:
   \[ \begin{pmatrix}
      1 & 2 & 3 & 4 \\
      2 & 4 & 3 & 1
      \end{pmatrix}
      \circ
      \begin{pmatrix}
      1 & 2 & 3 & 4 \\
      2 & 1 & 3 & 4
      \end{pmatrix}
      =
      \begin{pmatrix}
      1 & 2 & 3 & 4 \\
      4 & 2 & 3 & 1
      \end{pmatrix}
   \]
\end{example}

